% ============================================================================
% LACC: Nén ngữ cảnh có nhận thức ngôn ngữ cho LLM tiếng Việt
% Bản tham khảo tiếng Việt — ĐỒNG BỘ với lacc_icisn_en.tex (bản nộp ICISN)
% LƯU Ý: ICISN yêu cầu bài nộp bằng tiếng Anh; file này chỉ là tài liệu
% tham khảo/đối chiếu nội bộ, KHÔNG dùng để nộp hội nghị.
% ============================================================================
\documentclass[runningheads]{llncs}

\usepackage[T5,T1]{fontenc}
\usepackage[utf8]{inputenc}
\usepackage{vntex}
\usepackage{graphicx}
\usepackage{amsmath,amssymb}
\usepackage{booktabs}
\usepackage{array}
\usepackage{float}
\usepackage{hyperref}
\usepackage{color}
\renewcommand\UrlFont{\color{blue}\rmfamily}
\urlstyle{rm}

\newcommand{\tone}[1]{\textit{#1}}
\newcommand{\score}[1]{\mathcal{S}(#1)}
\newcommand{\weight}[1]{\mathbf{w}_{#1}}
\newcommand{\token}[1]{\texttt{#1}}
\newcommand{\method}[1]{\textsc{#1}}
\newcommand{\viet}[1]{\textit{#1}}

\begin{document}

\title{LACC: Nén Ngữ Cảnh Có Nhận Thức Ngôn Ngữ cho Mô Hình Ngôn Ngữ Lớn Tiếng Việt}

\titlerunning{LACC: Nén Ngữ Cảnh Có Nhận Thức Ngôn Ngữ cho LLM Tiếng Việt}

\author{Thân Ngọc Thiện}

\authorrunning{T. N. Thiện}

\institute{Trường Công nghệ Thông tin và Truyền thông,\\
Đại học Công nghiệp Hà Nội, Hà Nội, Việt Nam\\
\email{thien.than@haui.edu.vn}}

\maketitle

\begin{abstract}
Các phương pháp nén ngữ cảnh cho mô hình ngôn ngữ lớn (LLM) hiện nay — LLMLingua, SnapKV, H2O, StreamingLLM — đều \emph{bất chấp ngôn ngữ}: mọi token được xử lý như nhau bất kể ngôn ngữ nào. Với tiếng Việt — ngôn ngữ đơn lập, có 6 thanh điệu, trong đó 30--40\% token là hư từ — điều này gây ra mất thông tin thanh điệu khi nén, lãng phí ngân sách cho hư từ ít thông tin, và hiện tượng token inflation. Chúng tôi đề xuất \method{LACC} (\emph{Language-Aware Context Compression}), kết hợp một tín hiệu tương phản thanh điệu tính từ ma trận khoảng cách ngữ âm 6 thanh, một tín hiệu hình thái 5 lớp từ với ngân sách nén riêng cho từng lớp, và một tín hiệu perplexity, thành một điểm quan trọng token duy nhất, vận hành ở ba mức phần cứng từ 0~GB VRAM đến pipeline INT4 7B đầy đủ. Tín hiệu perplexity đến từ một mô hình nhỏ có sẵn hoặc từ một mô hình ngôn ngữ nhỏ (SLM) thuần Việt do chúng tôi huấn luyện với một hàm loss phụ trợ mới, \emph{Phonological Consistency Loss} (tối ưu đồng thời với LoRA), có bộ phân loại đóng vai trò kép làm tone probe lúc suy diễn. Chúng tôi báo cáo hiện trạng có thể kiểm chứng của hai nỗ lực đi kèm: \method{VCC-Bench} --- bộ benchmark nén ngữ cảnh tiếng Việt gồm 243 mẫu, 5 tác vụ, có quy trình xác thực nguồn gốc dữ liệu và tái lập kết quả; và một báo cáo huấn luyện/đánh giá từ đầu cho SLM tinh chỉnh LoRA này. Qua ba lần huấn luyện, việc mở rộng corpus từ 393 lên $\sim$22.200 đoạn văn đã nâng độ chính xác thanh điệu trên token có dấu thanh từ 16,6\% lên 92,9\%, và cải thiện perplexity validation 25,2\% so với mô hình gốc chưa tinh chỉnh. Một đánh giá thử nghiệm sơ bộ trên \method{VCC-Bench} cho thấy tín hiệu thanh điệu rule-based của \method{LACC} gần như tăng gấp đôi Tone Preservation Rate trung bình so với baseline không nén/ngẫu nhiên ở tỉ lệ nén 2$\times$ (0,96 so với 0,51--0,52). Chúng tôi báo cáo các kết quả này cùng các sửa lỗi kỹ thuật và vấn đề còn tồn đọng --- bao gồm một lỗi làm ROUGE-L mù trước dấu thanh tiếng Việt --- để các khẳng định trong bài đúng với những gì đã, và chưa, được đo đạc.

\keywords{Nén ngữ cảnh \and Mô hình ngôn ngữ lớn \and Xử lý ngôn ngữ tự nhiên tiếng Việt \and Bảo toàn thanh điệu \and Ngôn ngữ tài nguyên thấp \and Tinh chỉnh LoRA.}
\end{abstract}


\section{Giới thiệu}

Mô hình ngôn ngữ lớn ngày càng phải xử lý ngữ cảnh dài --- văn bản pháp luật, hội thoại nhiều lượt, chuỗi hành động agent nhiều bước --- trong khi chi phí $O(n^2)$ của self-attention khiến suy diễn đắt đỏ. Nén ngữ cảnh giải quyết vấn đề này theo ba hướng: \emph{nén prompt} (LLMLingua~\cite{jiang2023llmlingua}, Selective Context~\cite{li2023selective}) --- loại token có perplexity thấp; \emph{nén KV cache} (SnapKV~\cite{li2024snapkv}, H2O~\cite{zhang2023h2o}, StreamingLLM~\cite{xiao2024streamingllm}) --- loại bỏ vị trí cache theo điểm attention; và \emph{nén không gian tiềm ẩn} (Gist Tokens~\cite{mu2023gist}, ICAE~\cite{ge2024icae}) --- học memory token chuyên biệt.

Cả ba hướng đều \emph{mù ngôn ngữ}. Với tiếng Việt --- ngôn ngữ đơn lập, 6 thanh tương phản, khoảng 30--40\% token là hư từ --- điều này gây ra ba vấn đề cụ thể: \textbf{(P1) mất thông tin thanh điệu}, vì xóa token mang dấu thanh hoặc làm hỏng dấu có thể đổi hẳn nghĩa (\viet{má} ``mẹ'' $\to$ \viet{ma} ``con ma''); \textbf{(P2) lãng phí ngân sách cho hư từ}, vì các hư từ ít thông tin (\viet{đã}, \viet{của}, \viet{và}, \ldots) bị chấm điểm ngang hàng với thực từ; và \textbf{(P3) token inflation}, vì cùng nội dung tiếng Việt tốn 1.5--2.0$\times$ token so với tiếng Anh~\cite{lee2026equity,deng2026language}, làm hẹp context window hiệu dụng.

\textbf{Đóng góp.} (1) \method{LACC} --- khung tính điểm token có nhận thức ngôn ngữ, kết hợp tín hiệu tương phản thanh điệu, tín hiệu lớp hình thái, và tín hiệu perplexity thành một điểm quan trọng duy nhất, ở ba mức phần cứng (0~GB, $\sim$0.3--1~GB, INT4 7B đầy đủ). (2) \emph{Phonological Consistency Loss} --- một hàm mục tiêu huấn luyện tối ưu đồng thời với LoRA để hidden states của một SLM nhỏ mã hóa trực tiếp thanh điệu tiếng Việt; bộ phân loại thu được đóng vai trò kép làm \emph{tone probe} lúc suy diễn, khép kín vòng lặp huấn luyện--chấm điểm cho nhánh model-based (một đóng góp về phương pháp huấn luyện thực sự, khác với các tín hiệu dựa-từ-điển ở trên). (3) \method{VCC-Bench} --- bộ benchmark nén ngữ cảnh tiếng Việt đầu tiên, công bố cùng quy trình xác thực nguồn gốc/checksum/tái lập kết quả. (4) Một báo cáo huấn luyện và đánh giá từ đầu cho mô hình LoRA-cộng-tone-loss này, gồm ba lần huấn luyện, vấn đề quy mô dataset đã thúc đẩy các lần chạy này, và các vấn đề kỹ thuật --- kể cả một lỗi ROUGE-L nghiêm trọng --- hiện đang giới hạn những gì có thể khẳng định về chất lượng nén đầu-cuối. Chúng tôi tách bạch những gì đã \emph{thiết kế} (Mục~3), những gì đã \emph{đo được} (Mục~5), và những gì \emph{còn bỏ ngỏ} (Mục~6), để bài báo không phóng đại kết quả chưa được xác thực lại dưới một pipeline đã sửa lỗi.


\section{Nghiên cứu liên quan}

LLMLingua~\cite{jiang2023llmlingua} chấm điểm token bằng perplexity từ một LM tham chiếu nhỏ với budget controller, đạt tỉ lệ nén tới 20$\times$; Selective Context~\cite{li2023selective} chấm điểm bằng cosine similarity với câu hỏi. SnapKV~\cite{li2024snapkv} chọn vị trí KV cache từ attention trên một observation window ở cuối prompt; H2O~\cite{zhang2023h2o} giữ các ``heavy hitter'' có cumulative attention cao; StreamingLLM~\cite{xiao2024streamingllm} khai thác attention sink. Gist Tokens~\cite{mu2023gist} và ICAE~\cite{ge2024icae} thay vào đó huấn luyện mô hình nén prompt vào memory token đã học. Không phương pháp nào tích hợp tín hiệu đặc thù ngôn ngữ. Tiếng Việt là ngôn ngữ đơn lập, có thanh điệu với 6 thanh tương phản về ngữ âm~\cite{nguyen1997vietnamese}. Các nghiên cứu về công bằng tokenizer vẫn ghi nhận tiếng Việt chịu token inflation 1.5--2.0$\times$ so với tiếng Anh ngay cả với tokenizer công bằng nhất~\cite{lee2026equity}, và ngôn ngữ tài nguyên thấp chịu ``hình phạt kép'' về chi phí suy diễn cao hơn và độ chính xác thấp hơn~\cite{deng2026language}. Các mô hình tiếng Việt tiền huấn luyện như PhoBERT~\cite{nguyen2020phobert} cho thấy giá trị của mô hình hóa đặc thù ngôn ngữ cho NLP tiếng Việt nói chung, nhưng theo hiểu biết của chúng tôi, chưa có phương pháp công bố nào tích hợp tín hiệu thanh điệu hay hình thái vào chấm điểm quan trọng token cho bài toán nén ngữ cảnh, dù tiếng Việt hay bất kỳ ngôn ngữ có thanh điệu, tài nguyên thấp nào khác.


\section{Phương pháp đề xuất: LACC}

\method{LACC} tính điểm quan trọng cho mỗi token bằng cách kết hợp ba tín hiệu đã chuẩn hóa,
\begin{equation}
\score{t} = w_{\text{ppl}} \tilde{\score{}}_{\text{ppl}}(t) + w_{\text{tone}} \tilde{\score{}}_{\text{tone}}(t) + w_{\text{morph}} \tilde{\score{}}_{\text{morph}}(t),
\label{eq:combined}
\end{equation}
với mỗi $\tilde{\score{}}(\cdot)\!\in\![0,1]$ và tổng các trọng số bằng 1 ($w_{\text{ppl}}\!=\!0$, $w_{\text{tone}}\!=\!w_{\text{morph}}\!=\!0.5$ khi không dùng mô hình ngoài). $B=\max(\lfloor n/R\rfloor-2k,0)$ token ở giữa có điểm $\score{t}$ cao nhất được giữ lại, cùng $k{=}2$ token biên mỗi phía, theo đúng thứ tự gốc.

\textbf{Tín hiệu thanh điệu.} Với token đã giải mã $t$, gọi $\rho(t)$ là tỉ lệ ký tự mang thanh khác \tone{ngang} và $\nu(t)$ là số thanh khác-\tone{ngang} phân biệt xuất hiện. Trọng số bảo toàn là $\weight{tone}(t) = 1.0+\alpha\rho(t)(1+\beta\nu(t)/6)$, với $\alpha{=}0.5$, $\beta{=}0.3$ (khoảng $[1.0,1.65]$). Một ma trận tương phản ngữ âm $6\times6$, $D_{\text{tone}}(a,b)\in[0,0.9]$, đo mức độ khác biệt cảm nhận giữa hai thanh; với lân cận $\pm2$-token $N(t)$, hệ số tương phản là
\begin{equation}
f_{\text{contrast}}(t) = 1+\gamma\cdot\tfrac{1}{|N(t)|}\textstyle\sum_{n\in N(t)} D_{\text{tone}}(\text{tone}(t),\text{tone}(n)), \quad \gamma{=}0.4,
\end{equation}
và $\score{}_{\text{tone}}(t)=\weight{tone}(t)\times f_{\text{contrast}}(t)$. Ví dụ: \token{má} (\tone{sắc}) đứng cạnh \token{ma} (\tone{ngang}) và \token{mà} (\tone{huyền}) có điểm $\approx2.01$, gần gấp đôi điểm $\approx1.0$ của một hư từ trung tính về thanh điệu.

\textbf{Tín hiệu hình thái từ.} Mỗi token được tra từ điển và gán vào một trong 5 lớp --- \textsc{func} (hư từ, $\sim$35\%), \textsc{content} (thực từ, $\sim$50\%), \textsc{redup} (từ láy, $\sim$5\%), \textsc{compound} (từ ghép, $\sim$8\%), \textsc{sino} (từ Hán-Việt, $\sim$2\%) --- dùng từ điển đóng ($\sim$200 hư từ, $\sim$60 morpheme Hán-Việt) cho tra cứu $O(1)$. Mỗi lớp có hệ số bảo toàn ($0.40, 1.20, 0.60, 1.50, 1.50$ tương ứng) và tỉ lệ giữ lại $r_c$ ($0.30/0.85/0.50/0.95$ cho \textsc{func}/\textsc{content}/\textsc{redup}/\textsc{compound}) dùng để phân bổ ngân sách nén theo tỉ lệ, $B_c = B\cdot|T_c|r_c/\sum_{c'}|T_{c'}|r_{c'}$. Các cặp từ láy $(t_L,t_R)$ được phát hiện có điểm $t_L$ nhân $1.5$ và $t_R$ đặt lại $0.1$.

\textbf{Tín hiệu perplexity từ mô hình có sẵn.} Khi có GPU, một mô hình ngôn ngữ nhân quả nhỏ có sẵn (135M--1.5B tham số, vd SmolLM2-135M hay Qwen2.5-0.5B, $\sim$0.3--1.0~GB ở INT4) chấm điểm token theo độ bất ngờ, $\score{}_{\text{ppl}}(t_i)=-\log P(t_i\mid t_{<i})$, theo cửa sổ trượt 512 token --- cùng nguyên lý LLMLingua~\cite{jiang2023llmlingua} nhưng mô hình nhỏ hơn $\sim$50 lần. Bộ tính điểm này không cần huấn luyện gì thêm; đây là một mô hình có sẵn, cắm-là-chạy.

\textbf{Phương pháp huấn luyện đề xuất: Phonological Consistency Loss.} Thay vì chỉ dựa vào bộ tính điểm có sẵn ở trên, chúng tôi đề xuất huấn luyện một SLM thuần Việt sao cho hidden states của nó \emph{mã hóa trực tiếp thanh điệu}, bằng cách thêm một mục tiêu phụ trợ vào quá trình tinh chỉnh LoRA chuẩn,
\begin{equation}
\mathcal{L}_{\text{total}} = \mathcal{L}_{\text{LM}} + \lambda_{\text{tone}}\,\mathcal{L}_{\text{tone}}, \qquad
\mathcal{L}_{\text{tone}} = \frac1N\sum_{i=1}^{N} \text{CE}\big(\text{MLP}(h_i),\, y_i^{\text{tone}}\big),
\label{eq:phon}
\end{equation}
với $h_i$ là hidden state tại vị trí token \emph{đã nén} $i$, $y_i^{\text{tone}}\in\{0,\ldots,6\}$ là nhãn thanh điệu đúng (\tone{ngang} cộng 5 thanh có dấu), và $\text{MLP}$ là một bộ phân loại 3 lớp (Linear$\to$GELU$\to$LayerNorm$\to$Linear) được huấn luyện \emph{đồng thời} với LoRA adapter, không phải fit lại sau. Lúc suy diễn, chính bộ phân loại này được tái sử dụng làm \emph{tone probe}: điểm thanh điệu dựa-mô-hình cho token $t_i$ là $\score{}_{\text{tone-model}}(t_i) = 1.0 + \max_k \text{softmax}(\text{MLP}(h_i))_k$, tức token nào mô hình tự tin về thanh điệu thì được chấm điểm quan trọng hơn để giữ lại. Điều này khép kín vòng lặp giữa huấn luyện và suy diễn --- cùng một mục tiêu dạy mô hình bảo toàn thanh điệu cũng chính là tín hiệu dùng để chọn token --- khác với điểm thanh điệu dựa-từ-điển ở trên vốn không cần huấn luyện gì cả. Mục~5 báo cáo độ chính xác phân loại thanh điệu và perplexity của mô hình này trên tập hold-out qua ba lần huấn luyện; tone probe đã huấn luyện \textbf{chưa} thay thế bộ tính điểm có sẵn ở trên bên trong pipeline \method{LACC} đầu-cuối (Mục~6).

\textbf{Độ phức tạp.} Tín hiệu thanh điệu và hình thái chạy $O(n)$ trên CPU, 0~VRAM; mỗi bộ tính điểm perplexity (có sẵn hoặc đã huấn luyện) là $O(n\cdot d^2)$ trên GPU; chuẩn hóa và chọn Top-$B$ chạy $O(n\log n)$. Tầng không-mô-hình chỉ cần \textbf{0~GB VRAM}; tầng nhẹ chỉ cần thêm $\sim$0.3--1.0~GB ngoài mô hình sinh văn bản.


\section{VCC-Bench và Quy trình Đánh giá}

\method{VCC-Bench} v1 (snapshot \texttt{2026-07-06}) hiện có \textbf{243 mẫu trên 5 tác vụ}: Long-Document QA (160, văn bản pháp luật/báo chí, 2K--32K token), Multi-turn Conversation (48), Cross-lingual Compression (18), Needle-in-Haystack (9), và Agent Tool-Calling (8) --- nội dung lấy từ Wikipedia tiếng Việt (CC~BY-SA~4.0), văn bản pháp luật nhúng sẵn (phạm vi công cộng), và template tổng hợp, mỗi file có checksum SHA-256 để phát hiện sai lệch. Quy mô này nhỏ hơn và nghiêng về Long-Document QA nhiều hơn so với thiết kế mục tiêu ban đầu (500/300/400/250/200); mở rộng các tác vụ còn thiếu là việc tương lai. Một \texttt{ExperimentConfig} dùng chung cố định seed \texttt{42} (Python/NumPy/PyTorch), một split cố định duy nhất \texttt{"full"}, các tỉ lệ nén chuẩn $\{2\times,4\times,8\times\}$, và giải mã greedy, đồng thời lưu \texttt{config.json}/\texttt{environment.json} (git commit, phiên bản dependency) cho mỗi lần chạy để truy nguyên được.

\textbf{Một lỗi độ đo đáng báo cáo.} \texttt{compute\_rouge\_l} khởi tạo \texttt{RougeScorer} mà không truyền tokenizer; tokenizer mặc định của thư viện loại bỏ mọi ký tự ngoài \texttt{[a-z0-9]}, bao gồm \emph{mọi} nguyên âm mang dấu thanh tiếng Việt (U+1EA0--U+1EF9). Do đó \viet{bàn}, \viet{bán}, \viet{bạn} --- ba từ khác nhau --- đều thành \texttt{['b','n']}, khiến ROUGE-L chấm 1,0 cho một đáp án sai chỉ khác đáp án đúng ở dấu thanh. Vì ROUGE-L chiếm 40\% điểm chất lượng tổng hợp, lỗi này âm thầm khiến mọi kết quả chất lượng trước đó mù trước chính hiện tượng bài báo này nói đến. Chúng tôi đã sửa (2026-08-20) bằng \texttt{VietnameseRougeTokenizer} giữ nguyên dấu thanh, tương tự cách LongBench~\cite{bai2023longbench} chấm tiếng Trung bằng jieba. \textbf{Mọi kết quả tính trước bản sửa này đều không hợp lệ và cần chạy lại}, đây là lý do bảng kết quả chất lượng đầy đủ trên VCC-Bench chưa có trong bài báo này (Mục~6).


\section{Báo Cáo Huấn Luyện và Đánh Giá Đến Thời Điểm Hiện Tại}
\label{sec:report}

Chúng tôi tinh chỉnh \texttt{chronopt-research/vietnamese-gpt2-base} (125,6M tham số) bằng LoRA~\cite{hu2021lora} ($r{=}8$, $\alpha{=}16$; 1,18M tham số huấn luyện, 0,94\%), tối thiểu hóa causal-LM loss cộng $\lambda_{\text{tone}}{=}0.1$ lần Phonological Consistency Loss, trên một GPU NVIDIA GTX~1060 (6~GB), batch size 1, độ dài chuỗi 128, gradient accumulation 8, FP16.

\textbf{Run~1 (vì sao phải mở rộng corpus).} Chỉ dùng corpus gốc 393 đoạn (3 epoch, mặc định), đánh giá trên 40 văn bản hold-out, tone accuracy trên token có dấu thanh chỉ đạt 16,6--19,6\% --- không phân biệt được với mức ngẫu nhiên $\sim$20\% cho 5 lớp thanh khác-\tone{ngang} (perplexity 51,57). Chúng tôi coi Run~1 là một phép chẩn đoán, không phải kết quả: 393 đoạn là corpus quá nhỏ để học thanh điệu vượt quá ngẫu nhiên.

\textbf{Run 2--3 (mở rộng corpus).} Chúng tôi xây một corpus lớn hơn (UVW-2026, CC-BY-SA~4.0, + một tuyển tập thơ tiếng Việt, MIT) với split 90/10 cố định (seed 42). Run~2 dùng tỉ lệ mặc định 5.000 UVW-2026/556 thơ (5.556 đoạn); Run~3 mở rộng UVW-2026 gấp 10 lần ($\sim$22.200 đoạn tổng), giảm epoch từ 3 xuống 2 để giữ thời gian huấn luyện hợp lý. Bảng~\ref{tab:slm_train}--\ref{tab:slm_eval} báo cáo cấu hình và đánh giá hold-out, mỗi cột trên đúng split validation đã lưu của lần chạy đó.

\begin{table}[H]
\centering
\caption{Cấu hình huấn luyện SLM, Run 2--3 (mô hình gốc: \texttt{vietnamese-gpt2-base}, 125,6M).}
\label{tab:slm_train}
\begin{tabular}{lcc}
\toprule
\textbf{Tham số} & \textbf{Run 2} & \textbf{Run 3} \\
\midrule
Corpus (UVW-2026 / thơ, split 90/10) & 5.000 / 556 & 20.000 / $\sim$2.200 \\
Số văn bản train / validation & $\sim$5.000 / 555 & $\sim$20.000 / 2.218 \\
Epochs (batch 1, len 128, accum.\ 8) & 3 & 2 \\
\bottomrule
\end{tabular}
\end{table}

\begin{table}[H]
\centering
\caption{Đánh giá hold-out SLM (base và LoRA chỉ so sánh trong cùng split của một run).}
\label{tab:slm_eval}
\begin{tabular}{lcccc}
\toprule
\textbf{Độ đo} & \textbf{R2 Base} & \textbf{R2 LoRA} & \textbf{R3 Base} & \textbf{R3 LoRA} \\
\midrule
Perplexity $\downarrow$ & 159,81 & 124,26 & 148,10 & 110,76 \\
$\Delta$ so với base (cùng split) & \multicolumn{2}{c}{$-22,2\%$} & \multicolumn{2}{c}{$-25,2\%$} \\
Tone acc., token có dấu thanh $\uparrow$ & -- & \textbf{82,24\%} & -- & \textbf{92,94\%} \\
Baseline majority-class & -- & 44,93\% & -- & 46,15\% \\
\bottomrule
\end{tabular}
\end{table}

Tone accuracy trên token có dấu thanh tăng đơn điệu theo quy mô corpus --- 16,6--19,6\% (Run~1, 393 văn bản) $\to$ 82,24\% (Run~2, 5.556 văn bản) $\to$ 92,94\% (Run~3, $\sim$22.200 văn bản) --- luôn cao hơn hẳn cả mức ngẫu nhiên ($\sim$20\%) lẫn baseline majority-class ($\sim$45--46\%), và LoRA mang lại cải thiện perplexity có thể tái lập so với base đóng băng, xác nhận độc lập trên hai split ($-22,2\%$, $-25,2\%$). Lưu ý: \texttt{trained\_slm/final} bị ghi đè giữa các lần chạy, nên Run~2 và Run~3 dùng \emph{hai} split validation khác nhau --- chỉ so sánh base-với-LoRA trong cùng run là hợp lệ, và chênh lệch perplexity base (159,81 so với 148,10) phản ánh split khác nhau, không phải mô hình gốc (đã đóng băng) thay đổi; công cụ xây corpus ghép các bài thơ ngắn để thỏa ngưỡng lọc độ dài, có thể thổi phồng perplexity tuyệt đối không liên quan chất lượng mô hình; và số optimizer step cùng VRAM tối đa của Run~3 không được ghi lại, một khoảng trống công cụ đo cho lần chạy tiếp theo.

\textbf{Thử nghiệm sơ bộ trên VCC-Bench.} Tách biệt với phần trên, một pilot ba tầng sớm hơn (2026-07-07, trước cả cấu hình thống nhất ở trên lẫn bản sửa ROUGE-L) đã kiểm tra sơ bộ khả năng bảo toàn thanh điệu từ đầu đến cuối trên một tập con nhỏ (3--5 mẫu/tác vụ, chỉ tỉ lệ $\{2\times,4\times\}$), báo cáo Tone Preservation Rate (TPR) và tỉ lệ nén thực tế (CR) nhưng không có ROUGE-L/BLEU. Bảng~\ref{tab:pilot} lấy trung bình TPR trên cả 5 tác vụ; số liệu rule-based nhất quán (trong sai số) ở cả ba tầng phần cứng no-model, lightweight (+Qwen2.5-0.5B FP16), và full/INT4, nên chỉ hiển thị một giá trị đại diện mỗi phương pháp.

\begin{table}[H]
\centering
\caption{Thử nghiệm sơ bộ: TPR trung bình trên 5 tác vụ VCC-Bench (2026-07-07; tập con nhỏ).}
\label{tab:pilot}
\begin{tabular}{lcc}
\toprule
\textbf{Phương pháp} & \textbf{TPR @ 2$\times$} & \textbf{TPR @ 4$\times$} \\
\midrule
Không nén / baseline ngẫu nhiên & 0,51--0,52 & 0,29 \\
Chỉ morphology (tầng no-model) & 0,474 & 0,249 \\
Chỉ perplexity ngoài (tầng lightweight) & 0,545 & 0,316 \\
Model-based tone (tầng full/INT4)$^{*}$ & 0,609 & 0,394 \\
\textbf{Rule-based tone (\method{LACC})} & \textbf{0,960--0,963} & \textbf{0,538--0,541} \\
\textbf{Rule-based combined (tone+morph)} & \textbf{0,768--0,772} & \textbf{0,533--0,541} \\
\bottomrule
\end{tabular}
\end{table}
{\small $^{*}$Tỉ lệ nén thực tế thấp hơn mục tiêu (1,81$\times$/3,43$\times$), hoàn toàn do tác vụ needle-in-haystack, nơi bộ tính điểm cửa sổ trượt nén chưa đủ với ngữ cảnh rất dài ($\ge$32K token) trong pilot này --- một vấn đề kỹ thuật còn mở, không phải đặc tính chung của model-based scoring.}

Tín hiệu thanh điệu rule-based đơn lẻ gần như tăng gấp đôi TPR trung bình so với baseline không nén/ngẫu nhiên ở 2$\times$ (0,96 so với 0,51--0,52) và ở 4$\times$ (0,54 so với 0,29), với 0 VRAM thêm; morphology đơn lẻ không giúp ích cho TPR (nó tối ưu nội dung ngữ nghĩa, không phải ngữ âm), vì vậy \emph{combined} thấp hơn chỉ-tone dưới trọng số mặc định bằng nhau của Công thức~\ref{eq:combined}.

\textbf{Các mốc kỹ thuật.} Song song với các kết quả trên: một pipeline CI (lint $\to$ test $\to$ CPU smoke test $\to$ checksum dataset) chạy trên mọi push/PR, hoàn toàn CPU; \texttt{ExperimentConfig} thống nhất và lưu config/environment (Mục~4); một cài đặt Tone Preservation Rate duy nhất, chuẩn hóa, thay cho công thức tạm bợ của pilot tháng 7; tài liệu nguồn gốc dữ liệu và checksum đầy đủ; bản sửa ROUGE-L (Mục~4); một công cụ xây corpus huấn luyện lớn hơn và một tác vụ QA downstream thật từ UIT-ViQuAD2.0, cả hai đã cài đặt nhưng chưa có số liệu báo cáo; và một nghiên cứu đối chứng tone probe (frozen-base so với LoRA) để tách bạch tín hiệu thanh điệu LoRA thực sự thêm được bao nhiêu, chưa chạy xong.


\section{Thảo luận, Hạn chế và Hướng phát triển}

Kết quả pilot (Bảng~\ref{tab:pilot}) minh họa cụ thể cho vấn đề P1: nén mù ngôn ngữ loại token có dấu thanh với tỉ lệ gần như chọn ngẫu nhiên (TPR $\approx$0,51--0,52), trong khi tín hiệu thanh điệu bảo toàn ở mức $\approx$0,96. Với vấn đề P2, ngân sách theo tỉ lệ lớp chỉ giữ 30\% token hư từ so với 85\% token thực từ --- với văn bản 1.000 token (350/500/150 hư từ/thực từ/khác) nén xuống 250 token, một phương pháp mù ngôn ngữ giữ $\sim$87 hư từ trong khi \method{LACC} chỉ giữ $\sim$35, chuyển $\sim$52 token sang thực từ; đây là lập luận toán ngân sách, chưa phải hiệu ứng đầu-cuối đã đo trên VCC-Bench. Kiến trúc được thiết kế để mở rộng (ma trận thanh điệu cho các ngôn ngữ có thanh khác, vd tiếng Trung/Thái; bộ phân loại hình thái cho các ngôn ngữ đơn lập khác), nhưng chưa được kiểm chứng ngoài tiếng Việt.

Hạn chế chính, theo mức độ cấp thiết: (1) đánh giá chất lượng VCC-Bench đầy đủ (243 mẫu $\times\{2,4,8\}\times$, ROUGE-L/BLEU/EM) chưa chạy lại kể từ bản sửa ROUGE-L, nên được báo cáo là đang chờ, không bịa số liệu; (2) bộ tính điểm SLM thuần Việt đã huấn luyện và probe nhưng chưa thay thế vào pipeline \method{LACC} thật, nên chưa có hiệu ứng chất lượng đầu-cuối đã đo; (3) pilot (Bảng~\ref{tab:pilot}) quy mô nhỏ, đã cũ, dùng công thức TPR thô hơn độ đo canonical; (4) bộ tính điểm model-based nén chưa đủ với ngữ cảnh dài trong pilot (một lỗi còn mở); (5) tín hiệu thanh điệu/hình thái phụ thuộc từ điển đóng, sẽ phân loại sai từ mới hoặc từ mượn; (6) 243 mẫu VCC-Bench nghiêng về Long-Document QA, needle-in-haystack (9) và agent tool-calling (8) cần mở rộng; (7) thuật toán cần thấy toàn bộ ngữ cảnh trước khi nén, chưa có biến thể streaming. Ưu tiên ngắn hạn: chạy lại đánh giá VCC-Bench đầy đủ và ablation study dưới pipeline đã sửa lỗi; cắm tone probe SLM đã tinh chỉnh LoRA vào bộ nén thật và đo tác động lên chất lượng nén; hoàn thành nghiên cứu đối chứng tone probe; mở rộng các tác vụ còn thiếu mẫu; và sửa vấn đề nén chưa đủ của bộ tính điểm model-based với ngữ cảnh dài.


\section{Kết luận}

Chúng tôi trình bày \method{LACC}, một khung nén ngữ cảnh có nhận thức ngôn ngữ kết hợp tín hiệu tương phản thanh điệu, lớp hình thái, và perplexity tùy chọn từ mô hình ngoài để bảo vệ thông tin đặc thù tiếng Việt khi nén, từ 0~GB VRAM đến sinh văn bản INT4 7B đầy đủ. Chúng tôi báo cáo, đầy đủ, hiện trạng của hai nỗ lực đi kèm: \method{VCC-Bench} --- bộ benchmark nén tiếng Việt 243 mẫu, 5 tác vụ, có tài liệu nguồn gốc dữ liệu; và một báo cáo huấn luyện/đánh giá cho một bộ tính điểm SLM tiếng Việt có nhận thức thanh điệu, trong đó mở rộng corpus huấn luyện $\sim$56 lần (393$\to$22.200 văn bản) đã đưa độ chính xác trên token có dấu thanh từ mức ngẫu nhiên (16,6\%) lên 92,9\%, cùng cải thiện perplexity 22--25\% có thể tái lập so với base đóng băng. Chúng tôi cũng phát hiện và sửa một lỗi ROUGE-L khiến mọi đánh giá chất lượng trước đó mù trước dấu thanh tiếng Việt --- một phát hiện chúng tôi cho là quan trọng để báo cáo không kém một kết quả tích cực. Tín hiệu thanh điệu rule-based đơn lẻ gần như tăng gấp đôi Tone Preservation Rate trung bình so với baseline không nén/ngẫu nhiên trong một thử nghiệm sơ bộ (0,96 so với 0,51--0,52 ở 2$\times$). Bước tiếp theo là một đánh giá VCC-Bench quy mô đầy đủ đã sửa lỗi và việc cắm bộ tính điểm SLM đã huấn luyện vào pipeline nén thật.


\begin{credits}
\subsubsection{Lời cảm ơn.} Tác giả xin cảm ơn TS. Đặng Trọng Hợp (Trường Công nghệ Thông tin và Truyền thông, Đại học Công nghiệp Hà Nội) đã hướng dẫn thực hiện nghiên cứu này.
\subsubsection{Xung đột lợi ích.} Tác giả không có xung đột lợi ích nào liên quan đến nội dung bài báo này. Nghiên cứu này là nguyên bản và chưa từng được công bố hoặc gửi đăng ở nơi khác.
\end{credits}

\begin{thebibliography}{12}

\bibitem{jiang2023llmlingua}
Jiang, H., Wu, Q., Lin, C.-Y., Yang, Y., Qiu, L.: LLMLingua: Compressing Prompts for Accelerated Inference of Large Language Models. In: Proceedings of EMNLP (2023)

\bibitem{li2024snapkv}
Li, Y., Huang, Y., Yang, B., Venkitesh, B., Locatelli, A., Ye, H., Cai, T., Lewis, P., Chen, D.: SnapKV: LLM Knows What You are Looking for Before Generation. arXiv:2404.14469 (2024)

\bibitem{zhang2023h2o}
Zhang, Z., Sheng, Y., Zhou, T., Chen, T., Zheng, L., Cai, R., Song, Z., Tian, Y., Ré, C., Barrett, C., et al.: H2O: Heavy-Hitter Oracle for Efficient Generative Inference of Large Language Models. In: NeurIPS (2023)

\bibitem{xiao2024streamingllm}
Xiao, G., Tian, Y., Chen, B., Han, S., Lewis, M.: Efficient Streaming Language Models with Attention Sinks. In: ICLR (2024)

\bibitem{mu2023gist}
Mu, J., Li, X.L., Goodman, N.: Learning to Compress Prompts with Gist Tokens. In: NeurIPS (2023)

\bibitem{ge2024icae}
Ge, T., Hu, J., Wang, L., Wang, X., Chen, S.-Q., Wei, F.: In-context Autoencoder for Context Compression in a Large Language Model. In: ICLR (2024)

\bibitem{li2023selective}
Li, Y., Dong, B., Lin, C., Guerin, F.: Compressing Context to Enhance Inference Efficiency of Large Language Models. arXiv:2310.06201 (2023)

\bibitem{nguyen1997vietnamese}
Nguyễn, Đ.-H.: Vietnamese: Tiếng Việt Không Son Phấn. London Oriental and African Language Library, vol.~9. John Benjamins, Amsterdam (1997)

\bibitem{lee2026equity}
Lee, K., Qorib, M.R., Soegeng, A.I., Ng, H.T.: Equity with Efficiency: An Empirical Study of Tokenizers for Multilingual Large Language Models. arXiv:2606.15044 (2026)

\bibitem{deng2026language}
Deng, N., Shen, A., Feng, Y., Nwatu, J., Chung, J.-W., Chowdhury, M., Chen, Y., Mihalcea, R.: The Language-Energy Divide: Measuring Energy Costs of Multilingual LLM Inference. arXiv:2606.21869 (2026)

\bibitem{nguyen2020phobert}
Nguyen, D.Q., Nguyen, A.T.: PhoBERT: Pre-trained Language Models for Vietnamese. In: Findings of EMNLP (2020)

\bibitem{bai2023longbench}
Bai, Y., et al.: LongBench: A Bilingual, Multitask Benchmark for Long Context Understanding. arXiv:2308.14508 (2023)

\bibitem{hu2021lora}
Hu, E.J., Shen, Y., Wallis, P., Allen-Zhu, Z., Li, Y., Wang, S., Wang, L., Chen, W.: LoRA: Low-Rank Adaptation of Large Language Models. arXiv:2106.09685 (2021)

\end{thebibliography}

\end{document}
